\documentclass[lang=cn,newtx,10pt,scheme=chinese,nofont]{elegantbook}

% ============================================================
% 中文字体设置（nofont 模式下由用户自行定义）
% 注意：必须使用 XeLaTeX 或 LuaLaTeX 编译
% 正文字体：霞鹜臻楷（LXGW ZhenKai Slab），仅 Regular 一款
% ============================================================
% 主字体（正文）：霞鹜臻楷
\setCJKmainfont{LXGWZhenKaiSlabGB-Regular}[
  Path = ./,              % 字体文件与本文件同目录
  Extension = .ttf,       % 字体为 .ttf 文件
  AutoFakeBold = 2.5,     % 只有 Regular，自动生成伪粗体
  ItalicFont = LXGWZhenKaiSlabGB-Regular,  % 斜体仍用同一字体（正文不变斜）
  ItalicFeatures = {Path=./, Extension=.ttf}
]
% 无衬线字体（\sffamily）：用系统黑体
\setCJKsansfont{SimHei}
% 等宽字体（\ttfamily）：用系统仿宋
\setCJKmonofont{FangSong}
% 字族与常用命令（供 elegantbook 的 \heiti/\kaishu 等使用）
\setCJKfamilyfont{zhsong}{LXGWZhenKaiSlabGB-Regular}[Path=./,Extension=.ttf]
\setCJKfamilyfont{zhhei}{SimHei}
\setCJKfamilyfont{zhkai}{KaiTi}
\setCJKfamilyfont{zhfs}{FangSong}
\newcommand*{\songti}{\CJKfamily{zhsong}}
\newcommand*{\heiti}{\CJKfamily{zhhei}}
\newcommand*{\kaishu}{\CJKfamily{zhkai}}
\newcommand*{\fangsong}{\CJKfamily{zhfs}}

% ---------- 取消正文斜体 ----------
% 1) 西文：斜体也指向正体（Times 正体）
\setmainfont{TeXGyreTermesX}[
  UprightFont = *-Regular,
  BoldFont = *-Bold,
  ItalicFont = *-Regular,        % 斜体=正体，正文西文不倾斜
  BoldItalicFont = *-Bold,
  Extension = .otf,
  Scale = 1.0
]
% 2) 中文：\itshape / \emph 不再倾斜（保持正体）
\xeCJKsetup{AutoFakeSlant=false}
% 3) \emph 在中文下也改用正体而非斜体
\AtBeginDocument{%
  \renewcommand{\em}{\normalfont}%
  \renewcommand{\emph}[1]{{\normalfont #1}}%
}

\title{A版必修一}
\subtitle{高中数学讲义}
\author{JIE杰}
\institute{JIE杰}
\date{\today}
\version{0.1}

% 如果有封面和 logo，取消下面注释，并确保图片在 figure/ 下
 \cover{cover.jpg}
% \logo{logo-blue.png}

\begin{document}

\maketitle

\frontmatter
\tableofcontents

\mainmatter

% ============================================================
% 正文：按章用 \input 拼接
% 新增章节步骤：
%   1) 在 chapters\ 下新建片段文件，如 chapters\06-平面向量.tex
%   2) 在下方复制一组 \chapter + \input，改编号和文件名
%   3) 片段文件内不要写 \chapter、\documentclass、\begin{document}
% ============================================================

\chapter{集合与常用逻辑用语}
% ============================================================
% 片段文件：第 1 章  集合与常用逻辑用语
% 由主文件 讲义.tex 通过 \input 调用
% 注意：本文件不含 \chapter，章标题写在主文件里
% ============================================================

\section{集合的概念}
\subsection{元素与集合}
\begin{definition}[元素]
元素是构成集合的对象。
\begin{example}
    \begin{itemize}
        \item 数字 1、2、3 是集合 $\{1,2,3\}$ 的元素。
        \item 字母 a、b、c 是集合 $\{a,b,c\}$ 的元素。
    \end{itemize}
\end{example}
\end{definition}
\begin{axiom}[集合的确定性]
给定一个集合，任何一个对象要么是这个集合的元素，要么不是这个集合的元素，二者必居其一。集合与其元素之间的关系由 \textbf{ZF} 公理系统刻画（如外延公理：两个集合相等当且仅当它们含有完全相同的元素）。
\end{axiom}
\begin{definition}[集合]
集合是具有某种共同特征的对象的全体。
\end{definition}

\begin{definition}[元素与集合的关系]
若 $a$ 是集合 $A$ 的元素，记作 $a \in A$；若 $a$ 不是集合 $A$ 的元素，记作 $a \notin A$。
\end{definition}
\begin{definition}[常用数集]
\begin{itemize}
    \item 全体非负整数组成的集合称为自然数集，记作 $\mathbb{N}$；
    \item 全体正整数组成的集合称为正整数集，记作 $\mathbb{N}^*$ 或 $\mathbb{N}_+$；
    \item 全体整数组成的集合称为整数集，记作 $\mathbb{Z}$；
    \item 全体有理数组成的集合称为有理数集，记作 $\mathbb{Q}$；
    \item 全体实数组成的集合称为实数集，记作 $\mathbb{R}$。
\end{itemize}
\end{definition}
\subsection{集合的表示方法}
\begin{definition}[列举法]
把集合中的所有元素一一列举出来，并用花括号 $\{\ \}$ 括起来表示集合的方法叫做列举法。
\end{definition}
\begin{definition}[描述法]
用集合所含元素的共同特征表示集合的方法叫做描述法，其一般形式为 $\{x \mid p(x)\}$，其中 $x$ 为元素，$p(x)$ 为元素满足的条件。
\end{definition}
\begin{definition}[图示法]
用平面上封闭曲线的内部表示集合的方法叫做图示法（韦恩图）。
\end{definition}
\subsection{集合中元素的特性}
\begin{definition}[确定性、互异性、无序性]
集合中的元素具有如下三个特性：
\begin{itemize}
    \item 确定性：给定集合，任何一个对象是不是它的元素是确定的；
    \item 互异性：集合中的元素互不相同；
    \item 无序性：集合中的元素没有先后顺序。
\end{itemize}
\end{definition}

\section{集合间的基本关系}
\subsection{子集与真子集}
\begin{definition}[子集]
若集合 $A$ 中的任意一个元素都是集合 $B$ 的元素，则称 $A$ 是 $B$ 的子集，记作 $A \subseteq B$（或 $B \supseteq A$）。\\
规定：空集是任何集合的子集。
\end{definition}
\begin{theorem}[子集的性质]
对于任意集合 $A$、$B$、$C$：
\begin{itemize}
    \item 自反性：$A \subseteq A$；
    \item 反对称性：若 $A \subseteq B$ 且 $B \subseteq A$，则 $A=B$；
    \item 传递性：若 $A \subseteq B$ 且 $B \subseteq C$，则 $A \subseteq C$。
\end{itemize}
\end{theorem}
\begin{corollary}[子集个数]
含有 $n$ 个元素的集合共有 $2^n$ 个子集，其中真子集有 $2^n-1$ 个，非空真子集有 $2^n-2$ 个。
\end{corollary}
\begin{definition}[真子集]
若 $A \subseteq B$，且存在 $b \in B$ 但 $b \notin A$，则称 $A$ 是 $B$ 的真子集，记作 $A \subsetneq B$（或 $B \supsetneq A$）。
\end{definition}
\subsection{集合相等}
\begin{definition}[集合相等]
若 $A \subseteq B$ 且 $B \subseteq A$，则称集合 $A$ 与集合 $B$ 相等，记作 $A = B$。
\end{definition}
\subsection{空集}
\begin{definition}[空集]
不含任何元素的集合叫做空集，记作 $\varnothing$。
\end{definition}

\section{集合的基本运算}
\subsection{并集}
\begin{definition}[并集]
由所有属于集合 $A$ 或属于集合 $B$ 的元素组成的集合，叫做 $A$ 与 $B$ 的并集，记作 $A \cup B$，即 $A \cup B = \{x \mid x \in A \text{ 或 } x \in B\}$。
\end{definition}
\begin{lemma}[并集的性质]
对于任意集合 $A$、$B$：
\begin{itemize}
    \item 交换律：$A \cup B=B \cup A$；
    \item 幂等律：$A \cup A=A$；
    \item $A \cup \varnothing=A$，$A \subseteq A \cup B$。
\end{itemize}
\end{lemma}
\subsection{交集}
\begin{definition}[交集]
由所有既属于集合 $A$ 又属于集合 $B$ 的元素组成的集合，叫做 $A$ 与 $B$ 的交集，记作 $A \cap B$，即 $A \cap B = \{x \mid x \in A \text{ 且 } x \in B\}$。
\end{definition}
\begin{theorem}[并集与交集的运算律]
对于任意集合 $A$、$B$、$C$：
\begin{itemize}
    \item 交换律：$A \cap B=B \cap A$；
    \item 结合律：$(A \cup B) \cup C=A \cup (B \cup C)$，$(A \cap B) \cap C=A \cap (B \cap C)$；
    \item 分配律：$A \cap (B \cup C)=(A \cap B) \cup (A \cap C)$，$A \cup (B \cap C)=(A \cup B) \cap (A \cup C)$。
\end{itemize}
\end{theorem}
\subsection{补集}
\begin{definition}[全集与补集]
若集合 $U$ 含有我们所研究问题中涉及的所有元素，则称 $U$ 为全集；对于集合 $A \subseteq U$，由 $U$ 中不属于 $A$ 的所有元素组成的集合，叫做 $A$ 在 $U$ 中的补集，记作 $\complement_U A$，即 $\complement_U A = \{x \mid x \in U \text{ 且 } x \notin A\}$。
\end{definition}
\begin{theorem}[补集的性质（德摩根律）]
设全集为 $U$，$A$、$B \subseteq U$：
\begin{itemize}
    \item $A \cup \complement_U A=U$，$A \cap \complement_U A=\varnothing$；
    \item $\complement_U(\complement_U A)=A$；
    \item $\complement_U(A \cup B)=\complement_U A \cap \complement_U B$，$\complement_U(A \cap B)=\complement_U A \cup \complement_U B$。
\end{itemize}
\end{theorem}

\section{充分条件与必要条件}
\subsection{命题及其真假}
\begin{definition}[命题]
可以判断真假的陈述句叫做命题。判断为真的命题叫做真命题，判断为假的命题叫做假命题。
\end{definition}
\subsection{充分条件、必要条件}
\begin{definition}[充分条件与必要条件]
对于命题“若 $p$，则 $q$”，若该命题为真命题，即 $p \Rightarrow q$，则称 $p$ 是 $q$ 的充分条件，$q$ 是 $p$ 的必要条件。
\end{definition}
\subsection{充要条件}
\begin{definition}[充要条件]
若 $p \Rightarrow q$ 且 $q \Rightarrow p$，即 $p \Leftrightarrow q$，则称 $p$ 是 $q$ 的充分必要条件，简称充要条件。
\end{definition}
\begin{definition}[充分不必要、必要不充分、既不充分也不必要]
\begin{itemize}
    \item 若 $p \Rightarrow q$ 但 $q \nRightarrow p$，则 $p$ 是 $q$ 的充分不必要条件；
    \item 若 $q \Rightarrow p$ 但 $p \nRightarrow q$，则 $p$ 是 $q$ 的必要不充分条件；
    \item 若 $p \nRightarrow q$ 且 $q \nRightarrow p$，则 $p$ 是 $q$ 的既不充分也不必要条件。
\end{itemize}
\end{definition}

\section{全称量词与存在量词}
\subsection{全称量词与全称量词命题}
\begin{definition}[全称量词与全称量词命题]
“所有的”“任意一个”等短语叫做全称量词，用符号“$\forall$”表示；含有全称量词的命题叫做全称量词命题，一般形式为“$\forall x \in M,\ p(x)$”。
\end{definition}
\subsection{存在量词与存在量词命题}
\begin{definition}[存在量词与存在量词命题]
“存在一个”“至少有一个”等短语叫做存在量词，用符号“$\exists$”表示；含有存在量词的命题叫做存在量词命题，一般形式为“$\exists x \in M,\ p(x)$”。
\end{definition}
\subsection{全称量词命题与存在量词命题的否定}
\begin{definition}[全称量词命题的否定]
全称量词命题“$\forall x \in M,\ p(x)$”的否定是存在量词命题“$\exists x \in M,\ \neg p(x)$”。
\end{definition}
\begin{definition}[存在量词命题的否定]
存在量词命题“$\exists x \in M,\ p(x)$”的否定是全称量词命题“$\forall x \in M,\ \neg p(x)$”。
\end{definition}
\begin{theorem}[命题与其否定]
对任意命题 $p$，命题 $p$ 与其否定 $\neg p$ 必一真一假，不能同时为真也不能同时为假。
\end{theorem}
\begin{corollary}[全称量词命题为假的判定]
全称量词命题“$\forall x \in M,\ p(x)$”为假，当且仅当存在 $x_0 \in M$ 使 $p(x_0)$ 为假。
\end{corollary}


\chapter{一元二次函数、方程和不等式}
% ============================================================
% 片段文件：第 2 章  一元二次函数、方程和不等式
% 由主文件 讲义.tex 通过 \input 调用
% 注意：本文件不含 \chapter，章标题写在主文件里
% ============================================================

\section{等式性质与不等式性质}
\subsection{等式的基本性质}
\begin{definition}[等式的基本性质]
等式具有如下基本性质：
\begin{itemize}
    \item 对称性：若 $a=b$，则 $b=a$；
    \item 传递性：若 $a=b$，$b=c$，则 $a=c$；
    \item 可加性：若 $a=b$，则 $a+c=b+c$；
    \item 可乘性：若 $a=b$，则 $ac=bc$；
    \item 可除性：若 $a=b$ 且 $c \neq 0$，则 $\dfrac{a}{c}=\dfrac{b}{c}$。
\end{itemize}
\end{definition}
\begin{axiom}[实数的有序性]
任意两个实数 $a$、$b$，$a>b$、$a=b$、$a<b$ 三种关系中有且仅有一种成立。
\end{axiom}
\subsection{不等式的基本性质}
\begin{definition}[不等式的基本性质]
不等式具有如下基本性质：
\begin{itemize}
    \item 对称性：若 $a>b$，则 $b<a$；
    \item 传递性：若 $a>b$，$b>c$，则 $a>c$；
    \item 可加性：若 $a>b$，则 $a+c>b+c$；
    \item 可乘性：若 $a>b$，$c>0$，则 $ac>bc$；若 $a>b$，$c<0$，则 $ac<bc$；
    \item 同向可加性：若 $a>b$，$c>d$，则 $a+c>b+d$；
    \item 同向同正可乘性：若 $a>b>0$，$c>d>0$，则 $ac>bd$。
\end{itemize}
\end{definition}
\subsection{比较大小的方法}
\begin{definition}[作差法比较大小]
对于两个实数 $a$、$b$：若 $a-b>0$，则 $a>b$；若 $a-b=0$，则 $a=b$；若 $a-b<0$，则 $a<b$。
\end{definition}
\begin{theorem}[基本不等式的等价形式]
若 $a$、$b$ 为正数，则 $\dfrac{2}{\dfrac{1}{a}+\dfrac{1}{b}} \leq \sqrt{ab} \leq \dfrac{a+b}{2} \leq \sqrt{\dfrac{a^2+b^2}{2}}$，当且仅当 $a=b$ 时各等号均成立。
\end{theorem}
\begin{corollary}[重要不等式]
对任意实数 $a$、$b$，都有 $a^2+b^2 \geq 2ab$，当且仅当 $a=b$ 时等号成立。
\end{corollary}

\section{基本不等式}
\subsection{基本不等式及其几何意义}
\begin{definition}[基本不等式]
若 $a>0$，$b>0$，则 $\sqrt{ab} \leq \dfrac{a+b}{2}$，当且仅当 $a=b$ 时等号成立。其中 $\dfrac{a+b}{2}$ 叫做 $a$、$b$ 的算术平均数，$\sqrt{ab}$ 叫做 $a$、$b$ 的几何平均数。
\end{definition}
\subsection{基本不等式的应用}
\begin{definition}[基本不等式的常用变形]
由基本不等式可得：$a+b \geq 2\sqrt{ab}$（$a>0$，$b>0$），$ab \leq \left(\dfrac{a+b}{2}\right)^2$（$a$，$b \in \mathbb{R}$），当且仅当 $a=b$ 时等号成立。
\end{definition}

\section{二次函数与一元二次方程、不等式}
\subsection{一元二次方程的根与二次函数的图象}
\begin{definition}[一元二次函数]
形如 $y=ax^2+bx+c$（$a \neq 0$）的函数叫做一元二次函数，其图象是一条抛物线。
\end{definition}
\begin{definition}[一元二次方程的根]
使一元二次方程 $ax^2+bx+c=0$（$a \neq 0$）成立的实数 $x$ 叫做该方程的根，也叫做相应二次函数图象与 $x$ 轴交点的横坐标。
\end{definition}
\subsection{一元二次不等式及其解法}
\begin{definition}[一元二次不等式]
只含有一个未知数，并且未知数的最高次数是 $2$ 的不等式叫做一元二次不等式，其一般形式为 $ax^2+bx+c>0$ 或 $ax^2+bx+c<0$（$a \neq 0$）。
\end{definition}
\begin{theorem}[韦达定理]
设一元二次方程 $ax^2+bx+c=0$（$a \neq 0$，$\Delta \geq 0$）的两根为 $x_1$、$x_2$，则 $x_1+x_2=-\dfrac{b}{a}$，$x_1 x_2=\dfrac{c}{a}$。
\end{theorem}
\begin{corollary}[一元二次不等式的解集与判别式]
对于 $ax^2+bx+c>0$（$a>0$）：当 $\Delta>0$ 时解集为 $\{x \mid x<x_1 \text{ 或 } x>x_2\}$；当 $\Delta=0$ 时解集为 $\{x \mid x \neq -\tfrac{b}{2a}\}$；当 $\Delta<0$ 时解集为 $\mathbb{R}$。
\end{corollary}
\subsection{一元二次不等式的应用}
\begin{definition}[一元二次不等式的解法]
解一元二次不等式 $ax^2+bx+c>0$（$a>0$）：当 $\Delta>0$ 时，解集为 $\{x \mid x<x_1 \text{ 或 } x>x_2\}$；当 $\Delta=0$ 时，解集为 $\{x \mid x \neq x_1\}$；当 $\Delta<0$ 时，解集为 $\mathbb{R}$（其中 $x_1$，$x_2$ 为对应方程的根，$x_1<x_2$）。
\end{definition}


\chapter{函数的概念与性质}
% ============================================================
% 片段文件：第 3 章  函数的概念与性质
% 由主文件 讲义.tex 通过 \input 调用
% 注意：本文件不含 \chapter，章标题写在主文件里
% ============================================================

\section{函数的概念及其表示}
\subsection{函数的概念}
\begin{definition}[函数]
设 $A$、$B$ 是非空实数集，如果按照某种确定的对应关系 $f$，使对于集合 $A$ 中的任意一个数 $x$，在集合 $B$ 中都有唯一确定的数 $y$ 和它对应，那么就称 $f:A \to B$ 为从集合 $A$ 到集合 $B$ 的一个函数，记作 $y=f(x)$，$x \in A$。
\end{definition}
\begin{definition}[定义域与值域]
在函数 $y=f(x)$，$x \in A$ 中，$x$ 叫做自变量，$x$ 的取值范围 $A$ 叫做函数的定义域；与 $x$ 的值相对应的 $y$ 值叫做函数值，函数值的集合 $\{f(x) \mid x \in A\}$ 叫做函数的值域。
\end{definition}
\subsection{区间的概念}
\begin{definition}[区间]
设 $a$，$b$ 是两个实数，且 $a<b$：
\begin{itemize}
    \item 闭区间：$[a,b]=\{x \mid a \leq x \leq b\}$；
    \item 开区间：$(a,b)=\{x \mid a<x<b\}$；
    \item 半开半闭区间：$[a,b)=\{x \mid a \leq x<b\}$，$(a,b]=\{x \mid a<x \leq b\}$；
    \item 无穷区间：$[a,+\infty)=\{x \mid x \geq a\}$，$(-\infty,b]=\{x \mid x \leq b\}$ 等。
\end{itemize}
\end{definition}
\subsection{函数的表示法}
\begin{definition}[函数的表示法]
函数常用的表示方法有：解析法（用数学表达式表示）、列表法（用表格表示）和图象法（用图象表示）。
\end{definition}
\subsection{分段函数}
\begin{definition}[分段函数]
在定义域的不同子集上用不同解析式表示的函数叫做分段函数。
\end{definition}
\begin{axiom}[函数相等的判定]
两个函数相等的充要条件是它们的定义域相同，且对于定义域内任意 $x$，对应的函数值都相等。
\end{axiom}

\section{函数的基本性质}
\subsection{单调性与最大（小）值}
\begin{definition}[增函数与减函数]
设函数 $f(x)$ 的定义域为 $I$，区间 $D \subseteq I$。若对 $D$ 内任意 $x_1<x_2$，都有 $f(x_1)<f(x_2)$，则称 $f(x)$ 在 $D$ 上单调递增，$D$ 为其单调递增区间；若都有 $f(x_1)>f(x_2)$，则称 $f(x)$ 在 $D$ 上单调递减。
\end{definition}
\begin{theorem}[单调性的运算性质]
设 $f(x)$、$g(x)$ 在区间 $D$ 上单调递增，则 $f(x)+g(x)$ 在 $D$ 上单调递增；若 $f(x)$ 在 $D$ 上单调递增且 $f(x)>0$，则 $\dfrac{1}{f(x)}$ 在 $D$ 上单调递减。
\end{theorem}
\begin{definition}[最大值与最小值]
设函数 $y=f(x)$ 的定义域为 $I$。若存在 $x_0 \in I$，使得对任意 $x \in I$ 都有 $f(x) \leq f(x_0)$，则称 $f(x_0)$ 为 $f(x)$ 的最大值；若都有 $f(x) \geq f(x_0)$，则称 $f(x_0)$ 为最小值。
\end{definition}
\subsection{奇偶性}
\begin{definition}[奇函数与偶函数]
设函数 $f(x)$ 的定义域关于原点对称。若对定义域内任意 $x$ 都有 $f(-x)=-f(x)$，则 $f(x)$ 为奇函数；若都有 $f(-x)=f(x)$，则 $f(x)$ 为偶函数。
\end{definition}
\begin{theorem}[奇偶函数的图象特征]
奇函数的图象关于原点对称，偶函数的图象关于 $y$ 轴对称；若奇函数 $f(x)$ 在 $x=0$ 处有定义，则 $f(0)=0$。
\end{theorem}
\begin{corollary}[奇偶性的运算]
在公共定义域上：奇函数与奇函数之积为偶函数，偶函数与偶函数之积为偶函数，奇函数与偶函数之积为奇函数。
\end{corollary}
\subsection{函数性质的综合应用}
\begin{definition}[函数性质的综合]
函数的单调性、奇偶性、周期性等性质常常综合考查，研究函数时需结合定义域、图象等综合分析。
\end{definition}

\section{幂函数}
\subsection{幂函数的定义}
\begin{definition}[幂函数]
一般地，形如 $y=x^\alpha$（$\alpha$ 为常数）的函数叫做幂函数。
\end{definition}
\begin{theorem}[幂函数过定点]
$y=x^\alpha$ 的图象过点 $(1,1)$，当 $\alpha<0$ 时图象还过点 $(1,1)$ 附近且不过原点。
\end{theorem}
\begin{corollary}[幂函数 $y=x^\alpha$ 在 $(0,+\infty)$ 上的单调性]
当 $\alpha>0$ 时，$y=x^\alpha$ 在 $(0,+\infty)$ 上单调递增；当 $\alpha<0$ 时，$y=x^\alpha$ 在 $(0,+\infty)$ 上单调递减。
\end{corollary}
\subsection{常见幂函数的图象与性质}
\begin{definition}[常见幂函数的性质]
常见的幂函数 $y=x$、$y=x^2$、$y=x^3$、$y=x^{\frac{1}{2}}$、$y=x^{-1}$ 的图象均过点 $(1,1)$，其定义域、值域、单调性和奇偶性随指数 $\alpha$ 的不同而不同。
\end{definition}

\section{函数的应用}
\subsection{函数模型及其应用}
\begin{definition}[函数模型]
用函数刻画实际问题中变量之间依赖关系的数学表达式叫做函数模型。
\end{definition}
\subsection{用函数模型解决实际问题}
\begin{definition}[用函数模型解题]
用函数模型解决实际问题的一般步骤：审题、建模、求解、检验、作答。
\end{definition}


\chapter{指数函数与对数函数}
% ============================================================
% 片段文件：第 4 章  指数函数与对数函数
% 由主文件 讲义.tex 通过 \input 调用
% 注意：本文件不含 \chapter，章标题写在主文件里
% ============================================================

\section{指数}
\subsection{根式与分数指数幂}
\begin{definition}[根式]
若 $x^n=a$（$n>1$，$n \in \mathbb{N}^*$），则 $x$ 叫做 $a$ 的 $n$ 次方根。当 $n$ 为奇数时，$x=\sqrt[n]{a}$；当 $n$ 为偶数且 $a>0$ 时，$x=\pm\sqrt[n]{a}$。
\end{definition}
\begin{definition}[分数指数幂]
规定正数的正分数指数幂 $a^{\frac{m}{n}}=\sqrt[n]{a^m}$（$a>0$，$m$，$n \in \mathbb{N}^*$，$n>1$）；正数的负分数指数幂 $a^{-\frac{m}{n}}=\dfrac{1}{a^{\frac{m}{n}}}$；$0$ 的正分数指数幂等于 $0$。
\end{definition}
\subsection{指数幂的运算性质}
\begin{definition}[指数幂的运算性质]
设 $a>0$，$b>0$，$r$，$s \in \mathbb{Q}$：
\begin{itemize}
    \item $a^r \cdot a^s=a^{r+s}$；
    \item $(a^r)^s=a^{rs}$；
    \item $(ab)^r=a^r b^r$。
\end{itemize}
\end{definition}
\begin{theorem}[对数恒等式]
设 $a>0$，$a \neq 1$，$N>0$，则 $a^{\log_a N}=N$；又 $\log_a a^b=b$。
\end{theorem}
\begin{corollary}[对数换底推论]
$\log_a b \cdot \log_b a=1$，$\log_a b \cdot \log_b c \cdot \log_c a=1$（$a$、$b$、$c$ 均为正数且不等于 $1$）。
\end{corollary}

\section{指数函数}
\subsection{指数函数的概念}
\begin{definition}[指数函数]
形如 $y=a^x$（$a>0$ 且 $a \neq 1$）的函数叫做指数函数，其定义域为 $\mathbb{R}$。
\end{definition}
\subsection{指数函数的图象与性质}
\begin{definition}[指数函数的性质]
指数函数 $y=a^x$（$a>0$，$a \neq 1$）的图象恒过点 $(0,1)$。当 $a>1$ 时，函数在 $\mathbb{R}$ 上单调递增；当 $0<a<1$ 时，函数在 $\mathbb{R}$ 上单调递减。其值域为 $(0,+\infty)$。
\end{definition}
\subsection{指数函数的应用}
\begin{definition}[指数增长模型]
形如 $y=N(1+p)^x$（$N$ 为初始量，$p$ 为增长率）的函数模型常用于刻画指数增长问题。
\end{definition}

\section{对数}
\subsection{对数的概念}
\begin{definition}[对数]
若 $a^x=N$（$a>0$，$a \neq 1$），则数 $x$ 叫做以 $a$ 为底 $N$ 的对数，记作 $x=\log_a N$，其中 $a$ 叫做底数，$N$ 叫做真数。特别地，以 $10$ 为底的对数叫常用对数，记作 $\lg N$；以 $\mathrm{e}$ 为底的对数叫自然对数，记作 $\ln N$。
\end{definition}
\subsection{对数的运算性质}
\begin{definition}[对数的运算性质]
设 $a>0$，$a \neq 1$，$M>0$，$N>0$：
\begin{itemize}
    \item $\log_a(MN)=\log_a M+\log_a N$；
    \item $\log_a \dfrac{M}{N}=\log_a M-\log_a N$；
    \item $\log_a M^n=n\log_a M$（$n \in \mathbb{R}$）。
\end{itemize}
\end{definition}
\subsection{换底公式}
\begin{definition}[换底公式]
$\log_a b=\dfrac{\log_c b}{\log_c a}$（$a>0$，$a \neq 1$，$c>0$，$c \neq 1$，$b>0$）。特别地，$\log_a b=\dfrac{1}{\log_b a}$。
\end{definition}

\section{对数函数}
\subsection{对数函数的概念}
\begin{definition}[对数函数]
形如 $y=\log_a x$（$a>0$ 且 $a \neq 1$）的函数叫做对数函数，其定义域为 $(0,+\infty)$。
\end{definition}
\subsection{对数函数的图象与性质}
\begin{definition}[对数函数的性质]
对数函数 $y=\log_a x$（$a>0$，$a \neq 1$）的图象恒过点 $(1,0)$。当 $a>1$ 时，函数在 $(0,+\infty)$ 上单调递增；当 $0<a<1$ 时，函数在 $(0,+\infty)$ 上单调递减。其值域为 $\mathbb{R}$。
\end{definition}
\subsection{对数函数的应用}
\begin{definition}[对数函数的应用]
对数函数常用于解决与指数增长、溶液酸碱度（$\mathrm{pH}$）、声强级等相关的问题。
\end{definition}

\section{函数的零点与二分法}
\subsection{函数的零点及方程的解}
\begin{definition}[函数的零点]
对于函数 $y=f(x)$，使 $f(x)=0$ 的实数 $x$ 叫做函数 $y=f(x)$ 的零点。函数的零点就是方程 $f(x)=0$ 的根，也是函数图象与 $x$ 轴交点的横坐标。
\end{definition}
\subsection{用二分法求方程的近似解}
\begin{definition}[二分法]
对于在区间 $[a,b]$ 上图象连续不断且 $f(a)f(b)<0$ 的函数 $y=f(x)$，通过不断把零点所在区间一分为二，使区间端点逐步逼近零点，进而得到零点近似值的方法叫做二分法。
\end{definition}
\begin{theorem}[函数零点存在定理]
若函数 $y=f(x)$ 在闭区间 $[a,b]$ 上图象连续不断，且 $f(a)f(b)<0$，则存在 $c \in (a,b)$，使 $f(c)=0$，即函数在 $(a,b)$ 内至少有一个零点。
\end{theorem}
\begin{corollary}[零点个数与单调性]
若 $y=f(x)$ 在 $[a,b]$ 上单调且图象连续不断，又 $f(a)f(b)<0$，则 $f(x)$ 在 $(a,b)$ 内有且仅有一个零点。
\end{corollary}

\section{函数模型及其应用}
\subsection{不同函数模型的增长差异}
\begin{definition}[不同函数模型的增长差异]
当 $x$ 充分大时，增长速度由慢到快依次为：对数函数 $y=\log_a x$（$a>1$）、幂函数 $y=x^n$（$n>0$）、指数函数 $y=a^x$（$a>1$），即“指数爆炸、对数缓慢”。
\end{definition}
\subsection{函数模型的选择与应用}
\begin{definition}[函数模型的选择]
选择函数模型时，应根据实际问题的增长速度特征，选择与之相符的一次、二次、指数、对数或幂函数模型进行分析和预测。
\end{definition}


\chapter{三角函数}
% ============================================================
% 片段文件：第 5 章  三角函数
% 由主文件 讲义.tex 通过 \input 调用
% 注意：本文件不含 \chapter，章标题写在主文件里
% ============================================================

\section{任意角和弧度制}
\subsection{任意角}
\begin{definition}[任意角]
角可以看成平面内一条射线绕着端点从一个位置旋转到另一个位置所成的图形。按逆时针方向旋转形成的角叫做正角，按顺时针方向旋转形成的角叫做负角，没有旋转的角叫做零角。所有与角 $\alpha$ 终边相同的角构成集合 $\{\beta \mid \beta=\alpha+k\cdot 360^\circ,\ k \in \mathbb{Z}\}$。
\end{definition}
\subsection{弧度制}
\begin{definition}[弧度制]
长度等于半径长的圆弧所对的圆心角叫做 $1$ 弧度的角，记作 $1\ \mathrm{rad}$。用弧度作为单位来度量角的单位制叫做弧度制。
\end{definition}
\subsection{角度与弧度的换算}
\begin{definition}[角度与弧度的换算]
$180^\circ=\pi\ \mathrm{rad}$，即 $1^\circ=\dfrac{\pi}{180}\ \mathrm{rad}$，$1\ \mathrm{rad}=\left(\dfrac{180}{\pi}\right)^\circ$。弧长公式 $l=|\alpha| r$，扇形面积公式 $S=\dfrac{1}{2}lr=\dfrac{1}{2}|\alpha|r^2$。
\end{definition}

\section{三角函数的概念}
\subsection{三角函数的定义}
\begin{definition}[三角函数的定义]
设 $\alpha$ 是一个任意角，其终边上任意一点 $P(x,y)$（异于原点），$r=\sqrt{x^2+y^2}$，则定义 $\sin\alpha=\dfrac{y}{r}$，$\cos\alpha=\dfrac{x}{r}$，$\tan\alpha=\dfrac{y}{x}$。
\end{definition}
\subsection{三角函数值的符号}
\begin{definition}[三角函数值的符号]
三角函数值的符号由角的终边所在象限决定：第一象限全为正；第二象限正弦为正；第三象限正切为正；第四象限余弦为正。
\end{definition}
\subsection{同角三角函数的基本关系}
\begin{definition}[同角三角函数的基本关系]
$\sin^2\alpha+\cos^2\alpha=1$，$\tan\alpha=\dfrac{\sin\alpha}{\cos\alpha}$（$\alpha \neq k\pi+\dfrac{\pi}{2}$，$k \in \mathbb{Z}$）。
\end{definition}
\begin{theorem}[同角关系的常用变形]
$\sin\alpha+\cos\alpha$、$\sin\alpha-\cos\alpha$、$\sin\alpha\cos\alpha$ 三者知其一可求其余；$(\sin\alpha \pm \cos\alpha)^2=1 \pm 2\sin\alpha\cos\alpha$。
\end{theorem}
\begin{corollary}[三角函数值的符号与象限]
由三角函数值在各象限的符号可以确定角所在的象限，进而由同角关系求出其余三角函数值。
\end{corollary}

\section{三角函数的诱导公式}
\subsection{诱导公式的推导}
\begin{definition}[诱导公式]
利用角 $2k\pi+\alpha$、$\pi+\alpha$、$-\alpha$、$\pi-\alpha$、$\dfrac{\pi}{2}-\alpha$ 等与角 $\alpha$ 的终边关系，可推出各组诱导公式，其规律可概括为“奇变偶不变，符号看象限”。
\end{definition}
\subsection{诱导公式的应用}
\begin{definition}[诱导公式的应用]
诱导公式可用于将任意角的三角函数化为锐角的三角函数，进而求值、化简和证明。
\end{definition}

\section{三角函数的图象与性质}
\subsection{正弦函数、余弦函数的图象}
\begin{definition}[正弦函数、余弦函数的图象]
正弦函数 $y=\sin x$ 和余弦函数 $y=\cos x$ 的图象分别叫做正弦曲线和余弦曲线，可通过“五点法”作出一个周期内的简图。
\end{definition}
\subsection{正弦函数、余弦函数的性质}
\begin{definition}[正弦函数、余弦函数的性质]
$y=\sin x$ 与 $y=\cos x$ 的定义域均为 $\mathbb{R}$，值域均为 $[-1,1]$，最小正周期均为 $2\pi$。其中 $y=\sin x$ 为奇函数，在 $\left[-\dfrac{\pi}{2}+2k\pi,\dfrac{\pi}{2}+2k\pi\right]$ 上单调递增；$y=\cos x$ 为偶函数，在 $[-\pi+2k\pi,2k\pi]$ 上单调递增（$k \in \mathbb{Z}$）。
\end{definition}
\subsection{正切函数的图象与性质}
\begin{definition}[正切函数的图象与性质]
正切函数 $y=\tan x$ 的定义域为 $\{x \mid x \neq \dfrac{\pi}{2}+k\pi,\ k \in \mathbb{Z}\}$，值域为 $\mathbb{R}$，最小正周期为 $\pi$，是奇函数，在 $\left(-\dfrac{\pi}{2}+k\pi,\dfrac{\pi}{2}+k\pi\right)$ 上单调递增。
\end{definition}

\section{函数 \texorpdfstring{$y=A\sin(\omega x+\varphi)$}{y=Asin(ωx+φ)}}
\subsection{图象变换}
\begin{definition}[图象变换]
由 $y=\sin x$ 的图象经过振幅变换（$A$）、周期变换（$\omega$）和相位变换（$\varphi$）可得到 $y=A\sin(\omega x+\varphi)$ 的图象，其中 $A$ 为振幅，$\dfrac{2\pi}{\omega}$ 为周期，$\varphi$ 为初相。
\end{definition}
\subsection{性质与应用}
\begin{definition}[$y=A\sin(\omega x+\varphi)$ 的性质]
函数 $y=A\sin(\omega x+\varphi)$（$A>0$，$\omega>0$）的定义域为 $\mathbb{R}$，值域为 $[-A,A]$，最小正周期为 $\dfrac{2\pi}{\omega}$，常用于刻画简谐振动等周期现象。
\end{definition}
\begin{theorem}[图象的对称性]
$y=A\sin(\omega x+\varphi)$ 的对称轴方程为 $x=\dfrac{\pi}{2}-\varphi+k\pi$（$k \in \mathbb{Z}$），对称中心为 $\left(\dfrac{k\pi-\varphi}{\omega},0\right)$（$k \in \mathbb{Z}$）。
\end{theorem}
\begin{corollary}[五点法]
作出 $y=A\sin(\omega x+\varphi)$ 一个周期内的图象，可取 $\omega x+\varphi=0$、$\dfrac{\pi}{2}$、$\pi$、$\dfrac{3\pi}{2}$、$2\pi$ 所对应的五个关键点。
\end{corollary}

\section{三角恒等变换}
\subsection{两角和与差的正弦、余弦、正切公式}
\begin{definition}[两角和与差公式]
$\sin(\alpha \pm \beta)=\sin\alpha\cos\beta \pm \cos\alpha\sin\beta$；$\cos(\alpha \pm \beta)=\cos\alpha\cos\beta \mp \sin\alpha\sin\beta$；$\tan(\alpha \pm \beta)=\dfrac{\tan\alpha \pm \tan\beta}{1 \mp \tan\alpha\tan\beta}$。
\end{definition}
\subsection{二倍角公式}
\begin{definition}[二倍角公式]
$\sin 2\alpha=2\sin\alpha\cos\alpha$；$\cos 2\alpha=\cos^2\alpha-\sin^2\alpha=2\cos^2\alpha-1=1-2\sin^2\alpha$；$\tan 2\alpha=\dfrac{2\tan\alpha}{1-\tan^2\alpha}$。
\end{definition}
\begin{corollary}[降幂与半角公式]
$\sin^2\alpha=\dfrac{1-\cos 2\alpha}{2}$，$\cos^2\alpha=\dfrac{1+\cos 2\alpha}{2}$，$\sin\alpha\cos\alpha=\dfrac{1}{2}\sin 2\alpha$。
\end{corollary}
\subsection{简单的三角恒等变换}
\begin{definition}[辅助角公式]
$a\sin x+b\cos x=\sqrt{a^2+b^2}\sin(x+\varphi)$，其中 $\tan\varphi=\dfrac{b}{a}$，常用于化简三角函数式并求其最值。
\end{definition}
\begin{theorem}[正弦定理]
在 $\triangle ABC$ 中，$\dfrac{a}{\sin A}=\dfrac{b}{\sin B}=\dfrac{c}{\sin C}=2R$，其中 $R$ 为外接圆半径。
\end{theorem}
\begin{theorem}[余弦定理]
在 $\triangle ABC$ 中，$a^2=b^2+c^2-2bc\cos A$，$b^2=a^2+c^2-2ac\cos B$，$c^2=a^2+b^2-2ab\cos C$。
\end{theorem}
\begin{corollary}[三角形面积公式]
$S_{\triangle ABC}=\dfrac{1}{2}ab\sin C=\dfrac{1}{2}bc\sin A=\dfrac{1}{2}ac\sin B$。
\end{corollary}


%\chapter{平面向量}
%%#############################################################################
%  模板文件 —— 新建章节时请「复制本文件」再改名，不要直接在 \input 里写正文
%  本文件本身不被主文件调用，可以随意修改，不会影响编译结果
%#############################################################################
%
%  【一】放在哪里
%    chapters\<两位编号>-<章名>.tex
%    例如：chapters\06-平面向量.tex
%          chapters\07-复数.tex
%    编号决定排序，务必补零（01、02……10），不要写成 1、2、10。
%    章名建议与 \chapter{} 里的标题同名，便于查找。
%
%  【二】主文件里要加哪两行（在 讲义.tex 的 \mainmatter 之后）
%    \chapter{平面向量}
%    \input{chapters/06-平面向量}
%    注意：路径用正斜杠 / ，不要用反斜杠 \ ；扩展名 .tex 可写可不写。
%
%  【三】本文件里绝对不能出现的东西
%    \documentclass{...}      导言区已在主文件里，重复写会报错
%    \begin{document} ... \end{document}
%    \chapter{...}            章标题由主文件负责，这里只从 \section 开始
%    \usepackage{...}         宏包统一在主文件导言区添加
%    \setCJKmainfont 等字体命令  同上
%
%  【四】允许直接使用的东西
%    字体、颜色、定理环境、图片搜索路径（figure/、image/）都由主文件提供，
%    这里照常用即可。图片写成 \includegraphics{figure/xxx.png} 或直接写文件名。
%
%#############################################################################


% ======== 以下为示范内容，复制后请整体替换成你自己的正文 ========

\section{平面向量的概念}
\subsection{向量的定义}

\begin{definition}[向量]
既有大小又有方向的量叫做向量。
\end{definition}

% 可选的方括号里写「标题」，会显示为“定义 1.1（向量）”
\begin{definition}[零向量与单位向量]
长度为零的向量叫做零向量，记作 $\vec{0}$；长度等于 $1$ 的向量叫做单位向量。
\end{definition}

\begin{example}
下列各量中是向量的有（\quad）
\begin{enumerate}
    \item 质量
    \item 速度
    \item 温度
    \item 位移
\end{enumerate}
\end{example}

\subsection{向量的表示}

\begin{definition}[向量的几何表示]
用有向线段表示向量，有向线段的长度表示向量的大小，箭头所指的方向表示向量的方向。
\end{definition}

% 定理类环境都可带 [标题]，编号由主文件统一自动递增
\begin{theorem}[向量加法的三角形法则]
已知向量 $\vec{a}$、$\vec{b}$，在平面内任取一点 $A$，作 $\overrightarrow{AB}=\vec{a}$，$\overrightarrow{BC}=\vec{b}$，则 $\overrightarrow{AC}=\vec{a}+\vec{b}$。
\end{theorem}

\begin{lemma}[向量加法的交换律]
对任意向量 $\vec{a}$、$\vec{b}$，都有 $\vec{a}+\vec{b}=\vec{b}+\vec{a}$。
\end{lemma}

\begin{corollary}[向量不等式]
对任意向量 $\vec{a}$、$\vec{b}$，都有 $\bigl||\vec{a}|-|\vec{b}|\bigr| \leq |\vec{a}+\vec{b}| \leq |\vec{a}|+|\vec{b}|$。
\end{corollary}

% 需要插入图片时用 figure 环境；figure/ 和 image/ 已在主文件里设为自动搜索路径
% \begin{figure}[htbp]
%     \centering
%     \includegraphics[width=0.6\textwidth]{figure/示例图.png}
%     \caption{图的说明文字}
%     \label{fig:example}
% \end{figure}

\begin{note}
笔记环境没有编号，适合放易错提醒、补充说明。
\end{note}

\begin{remark}
注记环境用于补充评述，同样没有编号。
\end{remark}

% 带解答的练习题：exercise 自动编号，solution 不编号
\begin{exercise}
已知 $|\vec{a}|=3$，$|\vec{b}|=4$，$\vec{a} \perp \vec{b}$，求 $|\vec{a}+\vec{b}|$。
\end{exercise}
\begin{solution}
因为 $\vec{a} \perp \vec{b}$，所以 $\vec{a}\cdot\vec{b}=0$，故
\[
|\vec{a}+\vec{b}|=\sqrt{|\vec{a}|^2+2\vec{a}\cdot\vec{b}+|\vec{b}|^2}=\sqrt{9+0+16}=5.
\]
\end{solution}

% ======== 示范内容结束 ========


\end{document}